\documentclass[11pt,twoside,openright]{memoir}
\usepackage[T1]{fontenc}
\usepackage[utf8]{inputenc}
\usepackage{lmodern}
\usepackage{microtype}
\usepackage{graphicx}
\usepackage{xcolor}
\usepackage{booktabs}
\usepackage[hidelinks]{hyperref}
\setstocksize{9in}{6in}
\settrimmedsize{9in}{6in}{*}
\setlrmarginsandblock{0.82in}{0.68in}{*}
\setulmarginsandblock{0.82in}{0.75in}{*}
\checkandfixthelayout
\definecolor{chapterblue}{HTML}{24527A}
\chapterstyle{madsen}
\renewcommand{\chaptitlefont}{\normalfont\Huge\bfseries\color{chapterblue}}
\setsecnumdepth{subsection}
\begin{document}
\frontmatter
\begin{titlingpage}
\centering
\vspace*{1.1in}
{\Huge\bfseries A Field Guide to Quiet Machines\par}
\vspace{14pt}{\Large A multi-file book starter\par}
\vfill
{\Large Author Name\par}\vspace{8pt}
{\small A memoir-based workflow for books, essays, and print projects\par}
\end{titlingpage}
\cleartorecto
\chapter*{Preface}
This starter follows a gentle book-building workflow: establish the page,
separate front matter from chapters, edit in small files, and proof the final
sequence before sending it to print. The prose is intentionally replaceable.
\tableofcontents
\mainmatter
\chapter{Begin with the page}
\section{A useful first decision}
Every book project benefits from a small set of decisions made early: trim
size, running heads, chapter openings, and how illustrations will be supplied.
Those decisions become the stable frame around the writing.
\begin{quote}
The page is not a cage for the words. It is the instrument that lets their
rhythm be heard.
\end{quote}
\section{A repeatable structure}
Keep chapters independent while drafting. A chapter can be moved, split, or
revised without requiring the entire manuscript to be rewritten.
\chapter{Draft, proof, refine}
\section{A practical checklist}
\begin{enumerate}
  \item Read the text at the intended trim size.
  \item Check every figure, caption, cross-reference, and running head.
  \item Make a print proof before ordering a larger run.
\end{enumerate}
\section{A small production table}
\begin{center}
\begin{tabular}{llr}
\toprule Stage & Deliverable & Owner\\\midrule
Structure & chapter files and front matter & editor\\
Proofing & marked PDF and change list & author\\
Release & print PDF and source archive & production\\
\bottomrule
\end{tabular}
\end{center}
\backmatter
\chapter*{About this edition}
Replace this page with acknowledgements, colophon, edition history, or a final
call to action. The source is deliberately split into a single starter so it
can be expanded into one file per chapter as a project grows.
\end{document}
